\chapter{Examples, validation, and maintaining this manual}
\section{Included example inventory}
The following structures were generated through AtomForge's actual builder CLI, with the larger Cu host produced through the Python repetition API. They are reusable inputs for the walkthroughs.
\begin{longtable}{p{.34\linewidth}r p{.47\linewidth}}
\toprule File & Atoms & Construction \\ \midrule\endhead
\code{cu_fcc.cif} & 4 & Cubic 225, Cu at the origin, $a=3.61$ \angstrom. \\
\code{fe_bcc.cif} & 2 & Cubic 229, Fe at the origin, $a=2.87$ \angstrom. \\
\code{cu_host.cif} & 500 & $5\times5\times5$ repetition of the Cu conventional cell. \\
\code{cu_dislocation.cif} & 500 & Edge displacement, cylinder radius 5 \angstrom; family-change warning retained. \\
\code{cu_ni_alloy.cif} & 500 & Cu 0.7 / Ni 0.3, seed 42. \\
\code{cu_sphere.xyz} & 336 & Sphere radius 10 \angstrom, vacuum 5 \angstrom. \\
\code{cu_sigma5.cif} & 60 & Axis [001], $\Sigma5$, candidate 0, repetitions 3/3, overlap 1.5 \angstrom. \\
\code{cu_polycrystal.cif} & 1,329 & $25^3$ \angstrom$^3$, four grains, seed 7. \\
\code{sio2_pack.xyz} & 120 & Si 40 / O 80, density 1.5 $\mathrm{g\,cm^{-3}}$, seed 7. \\
\code{cu_mesh.xyz} & 224 & Unit octahedron OBJ at scale 12 \angstrom/model unit. \\
\bottomrule
\end{longtable}
\code{octahedron.obj} is an original procedural mesh included with the examples. The builder gallery plots actual output coordinates with Matplotlib. No external 3D model or Blender-generated asset is needed to reproduce these examples.

\section{CHGCAR provenance}
The original file remains at the user-supplied location and is not redistributed. \code{chgcar-summary.json} records its SHA-256 hash, dimensions, cell vectors, sites, units, range, integral, and periodicity. The manual includes derived \code{chgcar-planar.csv} and \code{chgcar-cumulative.csv} plus figures. A different CHGCAR will generally produce different numbers and images; the filename alone is not a unique identifier.

The example generation verifies that self-subtraction is exactly zero and that the final cumulative integral agrees with the full integral. These are numerical consistency checks, not independent validation of the originating electronic-structure calculation or a physical charge-transfer study.

Run \code{python docs/manual/scripts/check_examples.py} to verify the published structure counts and executable Python recipes. It uses an analytic periodic field to check charge conservation under complementary masks, Fourier reconstruction, derivative outputs, integration, neutral point-charge calls, and Cube/XSF round trips. This synthetic fixture is explicitly a software check and is separate from the supplied CHGCAR walkthrough.

\section{Regenerate examples and figures}
Build AtomForge and install the local Python package/native library first. The plotting script additionally needs NumPy and Matplotlib. Run from the repository root:
\begin{lstlisting}
python -m pip install numpy matplotlib
python docs/manual/scripts/generate_examples.py --chgcar "C:/path/to/CHGCAR"
\end{lstlisting}
Omit \code{--chgcar} to regenerate only structure examples and the builder gallery. Set \code{ATOMFORGE_PATH} if the executable is outside the normal build folder. The optional \code{MANUAL_PYTHON_DEPS} variable can point to a separate plotting-package directory when using an isolated embedded Python runtime.

The script overwrites the generated examples and scientific figures in this manual folder. Review new counts and the provenance JSON after changing software versions. It does not automate GUI screenshots: recapture those from the actual application when controls change. Keep scientific plots and interface captures labelled distinctly.

\section{Compile the LaTeX manual}
Install a LaTeX distribution providing pdfLaTeX and the packages used by \code{manual.tex}: Latin Modern, microtype, geometry, graphicx, xcolor, booktabs, longtable, tabularx, amsmath, enumitem, listings, fancyhdr, tcolorbox, and hyperref. Then run:
\begin{lstlisting}
python docs/manual/scripts/build_manual.py
\end{lstlisting}
The script compiles three passes for contents, page references, and PDF bookmarks. It places intermediates in \code{docs/manual/build/} and the distributable PDF at \code{docs/manual/AtomForge-manual.pdf}. The \code{build/} folder is ignored by Git. No raw CHGCAR is required to compile the checked-in manual because its published figures are included.

After editing, inspect the compiler log for overfull boxes, missing references, and missing characters. Render the PDF to page images and inspect figures, captions, tables, code wrapping, and page breaks. A successful LaTeX exit code alone does not establish a readable document. Keep screenshots at sufficient resolution for labels to remain legible at the chosen page width.

\section{Documentation coverage and limitations}
The manual covers the exposed File/Edit/Build/View/Analysis/Settings workflows, all 33 electronic operation choices, Python structure and electronic APIs, and CLI entry points. Optional stacking-fault controls are identified separately; unavailable interface/CNA/ADF menus are explicitly distinguished from exposed features. Screenshots cover representative builder and electronic workflows rather than every possible parameter combination.

The examples test executable generation and numerical consistency. They do not constitute convergence studies for every builder, material, or physical model. In particular, the single supplied CHGCAR does not establish a fragment-transfer reference set; random structures are not relaxed glasses; and model scattering densities are not self-consistent electronic densities.

\section{Source map for maintainers}
\begin{longtable}{p{.43\linewidth}p{.48\linewidth}}
\toprule Source area & Documentation to review when it changes \\ \midrule\endhead
\code{src/ui/FileBrowser.cpp} & Menu availability, file workflows, shortcuts, help and settings placement. \\
\code{src/ui/*BuilderDialog.cpp} & Builder controls, defaults, validation, preview/commit steps. \\
\code{src/algorithms/} & Builder geometry and structural-analysis assumptions. \\
\code{src/ui/ElectronicPostProcessing.cpp} & Tool labels, operand routing, imports/exports, calculation parameters. \\
\code{src/ui/ElectronicSliceViewport.cpp} & 2D controls, arbitrary-plane behaviour, preview limits. \\
\code{src/electronic/} & Units, endpoint conventions, numerical operations, Fourier and electrostatics. \\
\code{src/graphics/ElectronicViewport*} & Rendering, opacity, lighting, volume preview behaviour. \\
\code{python/atomforge/} & Python signatures, supported formats, library discovery, CLI choices. \\
\code{CMakeLists.txt}, CI workflows & Feature flags, dependencies, tests, packaging and interpreter selection. \\
\bottomrule
\end{longtable}
Update the source baseline on the title page when revising the manual. Re-run examples after algorithm or format changes, compare the generated metadata, and recapture affected UI figures. Avoid describing a feature as available solely because a corresponding source file exists.

\clearpage
\section{Glossary}
\begin{description}
\item[Asymmetric unit] Independent crystallographic sites expanded by space-group operations.
\item[Cell / origin] Lattice vectors and starting Cartesian position defining the grid or structure domain.
\item[Direct coordinates] Dimensionless coefficients of the direct lattice vectors.
\item[Field / channel] One scalar quantity sampled on a grid, such as total density or a derived difference.
\item[Reference / operand] The explicitly selected second input to an operation.
\item[Isosurface] A three-dimensional boundary at a constant scalar value.
\item[Contour] A constant-value line on a two-dimensional section.
\item[Volume rendering] Accumulation of contributions along viewing rays through a scalar field.
\item[Mask] Binary field defining a region within a grid domain.
\item[Reciprocal normal] A normal constructed from reciprocal lattice vectors; distinct from a direct-lattice direction in a general cell.
\item[Periodic endpoints] Repeated boundary samples that must be distinguished from unique periodic samples.
\item[Charge redistribution] Density change relative to specified references; not automatically a unique atom-assigned transfer.
\item[Voronoi partition] Geometric nearest-site regions, distinct from density-gradient basins.
\end{description}
